% 10_discussion.tex -- Section 9 Discussion, populated with empirical
% findings from the 2026-05-14 full-coverage real-data panel (12,895
% hourly Aave/Compound USDC overlapping observations, 98.5% joint
% coverage of 2024-11 to 2026-04).

\section{Discussion}
\label{sec:discussion}

The real-data panel assembled prior to test-window exposure
(\cref{sec:data}) surfaces three empirical features that bear directly
on the research questions of \cref{sec:methodology} and on the
architectural choices defended in \cite{solovev2026lob}. We discuss
each in turn, then respond head-on to the
Halpern--Pass--Saraf impossibility result.

\subsection{Governance Dynamics as a Structured Noise Source}
\label{sec:governance-noise}

A snapshot of the Aave V3 USDC reserve at 2026-05-14 returns a
post-kink slope $s_2 = 0.10$, sharply lower than the
historical $s_2 \approx 0.60$ that prevailed for most of the training
window. Cross-referencing on-chain
\texttt{ReserveInterestRateStrategyChanged} events places the change
in the Aave DAO's autumn 2025 risk-parameter cycle --- temporally
coincident with the 2025~Q3$\to$Q4 cross-protocol regime shift
documented in \cref{sec:regime-empirical}. The consequence for
forecasting is striking: subtracting
$f_{\text{kink}}(u;\theta_{\text{current}})$ from the realized supply
rate on the full 13{,}096-bar panel using
\emph{current} parameters explains only $21.3\%$ of supply-rate
variance, with residual mean $+186$\,bps and standard deviation
$227$\,bps --- far from the $60$--$80\%$ that a naive application of
the kink prior would suggest. Branch~A of the proposed forecaster,
whose target is precisely
$\varepsilon_i(t)=r_i(t)-f_{\text{kink}}(u_i(t);\theta_{\text{current}})$,
therefore absorbs a deterministic $+186$\,bps offset that would
otherwise be invisible to any architecture conditioning only on $r$
and $u$. The governance shift, identified in
\cref{sec:limitations} as a risk to the kink-subtraction design,
turns into a signal the forecaster can be expected to internalize.
A direct empirical prediction follows: ablation~7 (dual-branch
\emph{without} kink subtraction) should underperform the kinked
variant most clearly on the pre-governance-change portion of the
test window, and the gap should close on the post-change portion ---
a falsifiable architectural claim that the backtest will resolve.
Time-varying kink parameters obtained by streaming the full
\texttt{ReserveInterestRateStrategyChanged} event series into the
loader are the natural follow-up extension.

\subsection{Empirical Regime Structure Confirms Architectural Choice}
\label{sec:regime-empirical}

Quarterly aggregation of the cross-protocol spread on the
$12{,}895$-hour full-coverage overlap reveals a regime shift between
2025~Q3 and 2025~Q4 that is inconsistent with stationary dynamics.
The Aave-pays-more share moves from $39.9\%$ in Q3 (still
Compound-dominant; spread median $-0.218$\,pp) to $56.3\%$ in Q4
(Aave-dominant; spread median $+0.183$\,pp), a sixteen
percentage-point jump in conditional mean and the first sign flip of
the median spread on the panel --- the first quarter in which Aave
systematically pays more than Compound. Per-quarter spread standard
deviation across the panel ranges $0.57$--$5.75$\,pp; the Q3-to-Q4
shift in conditional mean is several multiples of the within-quarter
scale. The panel additionally exhibits a strong volatility regime
split: $\sigma=5.75$\,pp in 2024~Q4 (high-vol shock period) versus
$\sigma=0.57$\,pp in 2026~Q1 (calm), an approximately tenfold
variation that provides the natural anchor for the
tertile-conditional reporting of \cref{sec:regime}. These findings
constitute an empirical existence proof for the H2 architectural
claim: regime structure that single-state models cannot capture is
statistically present in the data, and a Markov-switching baseline
along the lines of \cite{hamilton1989regimeswitching} (ablation~4)
should show its largest gains over the EMA baseline precisely on the
Q3--Q4~2025 transition window. The dual-branch DA-BiGRU-CNN
architecture inherited from \cite{solovev2026lob} is designed for
exactly this kind of regime-dependent multivariate dynamics; the
Q3-to-Q4 shift makes the architectural choice empirically motivated
rather than merely transferred.

\subsection{Response to Halpern--Pass--Saraf on Short-Horizon Predictability}
\label{sec:fair-rates-response}

\cite{halpern2024fairrates} prove, by reduction of dynamic-repayment
lending to barrier-option pricing, that no \emph{fair} interest rate
--- in the sense of no-arbitrage under strategic borrowers who may
delay repayment --- exists for a realistic lending pool. They argue
on the strength of this result that time-series forecasting of pool
rates is methodologically misplaced. The $12{,}895$-hour
full-coverage cross-protocol panel sharpens the response. The
realized Aave-pays-more share is $40.7\%$, against $59.3\%$ for
Compound; the unconditional spread standard deviation is $2.51$\,pp,
with extremes $[-26.79, +44.89]$\,pp. A market in which cross-venue
rate inequality reverses on roughly a $41/59$ split, with material
per-hour dispersion and a documented intra-panel regime shift
(2025~Q3$\to$Q4), is by construction \emph{not} short-horizon
efficient: persistent allocation edge is available to a forecaster
that anticipates the crossovers ahead of the EMA-smoothed observer.
The Halpern--Pass--Saraf theorem operates at the equilibrium-existence
level and bounds the long-run fairness of any single rate; the
predictive MCDM operates at the trade-timing level and exploits
transient disequilibrium between two co-existing venues. The two
claims are orthogonal: the impossibility of a fair equilibrium rate
is fully compatible with, and in our panel empirically coexists with,
an exploitable short-horizon structure. The $41/59$ split with
$2.51$\,pp standard deviation is, in our reading, a directly measured
instance of the residual disequilibrium their result leaves
untouched.

\subsection{Limits of the LOB-to-DeFi Methodology Transfer}
\label{sec:transfer}

The transfer from \cite{solovev2026lob} succeeds at the level of the
dual-stream paradigm, the weighted-Pearson loss term, and the
ranking-oriented evaluation. It is re-derived, not mechanically
copied: LOB volume is unbounded and signed, while utilization is
bounded and clamped; the kink prior has no LOB analogue and
contributes the additional architectural lever that
\cref{sec:governance-noise} shows to carry real explanatory mass.
The principal residual risk is overfitting on a comparatively small
($\sim$7k bar) training set; we mitigate via dropout, walk-forward
cross-validation, and a deliberately compact ($\sim$50k parameter)
architecture. Under the null hypothesis H0
(\cref{sec:methodology}), the dual-branch architecture neither helps
nor hurts and the contribution reduces to the rigorous mainnet
backtest of MCDM-EMA, the falsification of the LOB-to-DeFi transfer
claim, and the open-source Extra+1/Extra+2 contributions of
\cref{sec:contributions-detail} --- each of which stands on its own
merits.

\subsection{Three Layers of Cross-Domain Transfer}
\label{sec:disc-three-layers}

The empirical record now spans two structurally different markets:
high-frequency LOB mid-price prediction \cite{solovev2026lob} and
the DeFi USDC lending-rate forecasting and allocation problem of
\cref{sec:defi-experiment}. The cross-domain hypothesis behind both
papers is not a single claim but a sequence of three: that an
identical \emph{architecture} carries signal across the two markets;
that the resulting \emph{predictions} carry the same kind of
out-of-sample structure; and that this in turn confers a
\emph{strategic} edge on an allocator built on top. The empirical
record now lets us disaggregate the three claims, and the conclusion
is that they fail at distinct layers --- the architectural transfer
succeeds, the predictive transfer is partial and regime-conditional,
and the strategic transfer fails on the binding test window. We
state each claim and the evidence in turn.

\subsubsection{Architectural Transfer Succeeds}
\label{sec:disc-arch}

The DA-BiGRU-CNN of \cref{sec:dual-branch} trains stably and
extracts non-trivial signal in both domains under \emph{identical}
network wiring, composite loss, and optimization schedule. In LOB
mid-price prediction it attains test weighted-Pearson on the order
of $0.74$, outperforming a 205-feature LightGBM baseline by
$0.5$--$1.6$ percentage points
\cite{solovev2026lob}; in DeFi lending it reaches validation
weighted-Pearson $0.747$ on Aave with directional accuracy $0.739$
and a positive Compound branch $R^{2}=0.138$
(\cref{sec:defi-forecaster-narrative}). The recipe that ports
without modification is concrete and recordable: (i) identify a
domain-natural decomposition pair $(\mathcal X_A, \mathcal X_B)$
such that one subspace carries the deterministic backbone and the
other carries the residual stochastic component --- (price, volume)
in LOB, $(\varepsilon = r-f_{\text{kink}}(u),\; u+\text{ctx})$ in
DeFi (\cref{sec:dual-branch}); (ii) re-use the same dual-BiGRU
stack with multi-scale Conv1d fusion (kernels $\{3,5,7\}$, hidden
size $64$); and (iii) re-use the same composite loss
(\cref{eq:loss}) with weighted-Pearson as the dominant term. The
\emph{architectural} transfer hypothesis --- ``the same template
generalises by re-identifying the decomposition pair'' --- is
therefore supported by the validation diagnostics on both panels.
This is the central methodological contribution of the present
paper and is independent of the H1 outcome on the strategic layer.

A subsidiary but important architectural finding is that the kink
prior has no LOB analogue and contributes a genuinely new lever in
the DeFi setting. The governance-noise analysis of
\cref{sec:governance-noise} shows that subtracting
$f_{\text{kink}}(u;\theta_{\text{current}})$ from the realized rate
explains only $21.3\%$ of supply-rate variance under static current
parameters; conditioning on the protocol-published kink function
absorbs a deterministic $+186$\,bps offset that would otherwise be
invisible to any architecture conditioning only on $(r, u)$. The
LOB-to-DeFi transfer is in this sense \emph{strictly enriching} at
the architectural layer: every LOB-side architectural choice ports,
and one DeFi-native lever (the kink prior) is added.

\subsubsection{Predictive Transfer is Partial}
\label{sec:disc-pred}

The transfer breaks one layer up. Validation weighted-Pearson
$0.747$ on the Aave branch collapses to test
$-0.07$, and directional accuracy on the cross-protocol label
falls from $0.739$ on validation to $0.351$ on test
(\cref{sec:defi-forecaster-narrative}) --- a sub-chance value
whose sign is the empirical signature of a forecaster that has
locked onto a correlation pattern from the wrong regime rather
than of a forecaster that is uncertain. The proximate cause is
documented and dated: between 2025-Q3 and 2025-Q4 the
Aave-pays-more share rises from $39.9\%$ to $56.3\%$ and the
median cross-protocol spread inverts from $-0.218$\,pp to
$+0.183$\,pp (\cref{sec:regime-empirical}); the validation fold
lives in the pre-inversion regime; the test fold straddles the
inversion and a return to Compound dominance in 2026-Q1. The
empirical contrast with the LOB benchmark of
\cite{solovev2026lob} is sharp. LOB microstructural dynamics on
millisecond horizons are dominated by stationary adverse-selection
mechanics whose statistical character does not shift on a
four-month time scale; DeFi lending rates inherit
on-the-order-of-quarter regime persistence from slow demand-side
composition shifts (USDC borrow-demand cohort changes, perpetual
DEX funding-rate cycles, periodic governance-set risk-parameter
updates), and on the panel assembled here a single quarter is the
characteristic regime time-scale rather than a sub-component of
within-regime noise. Predictive transfer is therefore conditional
on regime alignment, not architecturally guaranteed: in-distribution
predictability does not imply robust short-horizon predictability
in DeFi the way it appears to in microstructure.

A constructive reading of this partial-transfer result is the
quarterly-conditional retraining schedule of \cref{sec:future}.
The same architecture, trained on a rolling pre-inversion window
and re-fit at the Q3-to-Q4 boundary, is the obvious next
experiment. The negative test-fold weighted-Pearson does not
falsify the architecture; it falsifies the implicit assumption
that a one-shot training pass on the pre-inversion regime gives a
forecaster whose calibration survives the inversion.

\subsubsection{Strategic Transfer Fails on This Window}
\label{sec:disc-strat}

The predictive MCDM under-performs the reactive MCDM-EMA by
$8$\,bp in annualised net APY ($\PredictiveAPY{}$ vs
$\EMAAPY{}$); the paired monthly-Sharpe bootstrap of
\cref{sec:defi-h1-narrative} returns $\SharpeDelta{}$ with $95\%$
CI $[\BootstrapCILow{},\,\BootstrapCIHigh{}]$ and
$p=\BootstrapPvalue{}$, firmly inside the H0 acceptance region of
the pre-registered $\Delta\textsc{Sharpe}\ge 0.2$ test. The
hysteresis-sweep ablation of \cref{sec:defi-ablation-narrative}
rules out the most natural rebuttal --- ``the threshold is
mis-tuned'' --- by showing that the design $\theta=0.05$ sits at
the Sharpe knee of the sweep, with fewer rebalances (and zero
gas at $\theta\in\{0.10,0.20\}$) actually \emph{not raising} net
APY against EMA in this window. The mechanism is therefore: a
forecaster whose test-fold directional sign is wrong cannot drive
a switching allocator into the right pool ahead of the
EMA-smoothed reactive comparator; the predictive overhead consumes
gas without obtaining the timing alpha that justifies it.

This is the same negative-ensemble pattern that
\cite{solovev2026lob}~\S5.5 reported for the LOB setting ---
combining a strong dual-branch GRU forecaster with a 205-feature
LightGBM ranker \emph{degraded} test weighted-Pearson by $0.05$ to
$1.6$ percentage points rather than improving it --- generalised
one layer up the stack: combining a forecaster with a reactive
allocator can under-perform the allocator alone if the forecaster
transports out-of-distribution badly. Two negative results, in two
structurally different domains, with the same mechanistic
explanation: \emph{the apparent complementarity of two components
is a property of the joint distribution they were calibrated on,
not a robust architectural fact}. This is, in our reading, a
non-trivial finding that the H0-publishable protocol of
\cref{sec:contributions} is designed to surface. A literature in
which only H1-positive ensembles are reported is a literature in
which this cross-domain regularity is invisible.

\subsubsection{Relation to Halpern--Pass--Saraf, Revisited}
\label{sec:disc-impossibility}

The three-layer disaggregation lets us refine the response to
\cite{halpern2024fairrates} previewed in \cref{sec:fair-rates-response}.
Their result is a no-arbitrage statement about equilibrium fairness
at the contract-pricing layer; it constrains \emph{which interest
rate the protocol could in principle charge} under strategic
borrowers. Our finding is at the execution layer and it is
empirical rather than theoretical: short-horizon local
predictability of the realized rate stream is present
in-distribution (validation weighted-Pearson $0.747$ on Aave) but
does not survive the 2025-Q3-to-Q4 regime shift documented in the
data panel. These two negative findings do not refute each other ---
one is a no-arbitrage statement about equilibrium fairness, the
other is an out-of-sample statement about the stationarity of the
short-horizon predictive structure --- and they do not jointly
imply that DeFi-lending forecasting is hopeless. They jointly
imply that the empirical-predictability conjecture must be
re-stated in a regime-conditional form, with explicit handling of
distribution shift on the order of one fiscal quarter rather than
a global predictability claim across an 18-month panel. The
H0-publishable verdict is in this sense complementary to
\cite{halpern2024fairrates}: it provides the empirical data point
that the impossibility-of-fair-pricing result predicts will
manifest as regime-persistent disequilibrium in the realised rate
stream.
